% =============================================================================
% Training, Evaluation, and Results sections
% Requires: booktabs, graphicx, multirow, subcaption, hyperref
% Figures assumed in figures/slides/
% =============================================================================

% -----------------------------------------------------------------------------
\section{Training}
\label{sec:training}
% -----------------------------------------------------------------------------

We conduct three groups of pre-training experiments, each designed to isolate a
different axis of variation: model capacity, data quality at fixed compute, and
the interaction between data source and model architecture at full scale.
All runs use the Nemotron-Nano-V3 architecture family and are trained on
8$\times$NVIDIA H100 80\,GB GPUs.

\subsection{Architecture}
\label{sec:architecture}

\paragraph{Dense model.}
Our dense baseline follows the Nemotron-Nano-V3 design: a decoder-only
Transformer with Grouped-Query Attention (GQA), RMSNorm, SwiGLU activations,
and Rotary Position Embeddings (RoPE).  The 382M-parameter configuration uses
a hidden dimension of 1024, 24 layers, 16 attention heads (4 key-value heads),
an intermediate size of 4096, and a context length of 2048 tokens.  Input and
output embeddings are tied.

\paragraph{Mixture-of-Experts (MoE) model.}
The MoE variant replaces each dense feed-forward block with a top-2 gated
mixture of 8 experts, yielding 1.4B total parameters but only 390M active
parameters per token---roughly matching the dense model's per-token compute.
We use a hidden dimension of 768, 24 layers, 12 attention heads (4 key-value
heads), and an intermediate size of 3072 per expert.  A load-balancing
auxiliary loss ($\alpha=0.01$) encourages uniform expert utilization.

\paragraph{Tokenizer.}
All experiments share a 16K-vocabulary BPE tokenizer trained on the hybrid
pre-training corpus.

\subsection{Training Infrastructure}
\label{sec:infra}

Models are distributed across 8 GPUs using Fully Sharded Data Parallelism
(FSDP2).  We apply \texttt{torch.compile} for graph-level kernel fusion.
Training uses bfloat16 mixed precision throughout.

\paragraph{Optimizer.}
AdamW with $\beta_1=0.9$, $\beta_2=0.95$, $\epsilon=10^{-8}$, and weight
decay $0.1$ (applied to all parameters except biases and layer norms).
Gradients are clipped at a max norm of 1.0.

\paragraph{Learning rate schedule.}
Linear warmup over the first 1\% of training steps followed by cosine decay
to $0.5\times$ the peak learning rate of $3\times10^{-4}$.

\paragraph{Data loading.}
Large datasets (raw, hybrid) are stored as pre-sharded JSONL files and
consumed via a streaming data loader that tokenizes on-the-fly, avoiding the
need to materialize the full tokenized corpus in memory.  Smaller datasets
(curated) are tokenized in memory at startup.

\subsection{Experiment Group 1: Model Scaling}
\label{sec:group1}

The first group of experiments holds the data source constant (raw Wikipedia
and web text) and varies model capacity from 4M to 58M parameters
(Table~\ref{tab:group1}).  These runs establish a baseline scaling trend and
confirm that our training pipeline produces expected improvements with
increased capacity.

\begin{table}[h]
\centering
\caption{Group 1: Model scaling experiments (raw data, dense architecture).}
\label{tab:group1}
\begin{tabular}{lccc}
\toprule
\textbf{Experiment} & \textbf{Params} & \textbf{Tokens} & \textbf{Steps} \\
\midrule
\texttt{baseline\_29k\_docs} & 4M   & 26M   & 597    \\
\texttt{scaled\_38M}         & 39M  & 378M  & 720    \\
\texttt{raw\_sub\_57M}       & 58M  & 328M  & 396    \\
\texttt{curated\_57M}        & 58M  & 328M  & 936    \\
\texttt{hybrid\_sub\_57M}    & 58M  & 328M  & 936    \\
\texttt{scaled\_57M}         & 58M  & 378M  & 1,080  \\
\texttt{raw\_full\_57M}      & 58M  & 5.7B  & 2,525  \\
\bottomrule
\end{tabular}
\end{table}

We additionally vary data source at the 58M scale (raw, curated, hybrid,
scaled subsets) to provide an early signal on data quality effects before
committing to the larger runs in Groups 2 and 3.

\subsection{Experiment Group 2: Data Quality at Equal Token Budget}
\label{sec:group2}

To isolate the effect of data quality from data quantity, we train three
382M-parameter dense models on exactly 5.7B tokens each, drawn from
different corpora (Table~\ref{tab:group2}).  The curated corpus uses the
full available curated data ($\approx$5.3B tokens); for raw and hybrid we
subsample to match.

\begin{table}[h]
\centering
\caption{Group 2: Equal-token data quality ablation (382M dense, 5.7B tokens each).}
\label{tab:group2}
\begin{tabular}{lccc}
\toprule
\textbf{Experiment} & \textbf{Data Source} & \textbf{Tokens} & \textbf{Steps} \\
\midrule
\texttt{equal5b\_raw\_dense}    & Raw     & 5.7B & 10,871 \\
\texttt{full\_curated\_dense}   & Curated & 5.3B & 10,000 \\
\texttt{equal5b\_hybrid\_dense} & Hybrid  & 5.7B & 10,871 \\
\bottomrule
\end{tabular}
\end{table}

\paragraph{Data sources.}
\begin{itemize}
    \item \textbf{Raw}: A 25M-document mixture of FineWeb-Edu (17.5M),
          FinePDFs-Edu (5M), and FineWiki (2.5M), totaling $\approx$55B
          tokens.
    \item \textbf{Curated}: Quality-filtered subset of the raw corpus
          produced by perplexity filtering, deduplication, and domain
          balancing, yielding $\approx$5.3B tokens.
    \item \textbf{Hybrid}: The curated corpus augmented with
          LLM-generated synthetic data (via NVIDIA Data Curator),
          totaling $\approx$8.8B tokens.
\end{itemize}

\subsection{Experiment Group 3: Full-Scale Data $\times$ Architecture}
\label{sec:group3}

The final group trains both dense and MoE architectures on the full extent
of each data source (Table~\ref{tab:group3}), producing a $3\times2$ grid
of experiments that reveals how data quality and model architecture interact
at scale.  All models train for one full epoch over their respective corpus.

\begin{table}[h]
\centering
\caption{Group 3: Full-scale runs across data sources and architectures.}
\label{tab:group3}
\begin{tabular}{llcccc}
\toprule
\textbf{Data} & \textbf{Architecture} & \textbf{Total Params} & \textbf{Active Params} & \textbf{Tokens} & \textbf{Steps} \\
\midrule
Raw     & Dense & 382M  & 382M & 55.1B & 105,000 \\
Raw     & MoE   & 1.4B  & 390M & 55.1B & 210,000 \\
Curated & Dense & 382M  & 382M & 5.2B  & 10,000  \\
Curated & MoE   & 1.4B  & 390M & 6.6B  & 25,000  \\
Hybrid  & Dense & 382M  & 382M & 5.2B  & 10,000  \\
Hybrid  & MoE   & 1.4B  & 390M & 6.6B  & 25,000  \\
\bottomrule
\end{tabular}
\end{table}

The MoE models use a global batch size of 128 (16 per GPU $\times$ 8 GPUs)
while dense models use 256 (16 per GPU $\times$ 2 gradient accumulation
steps $\times$ 8 GPUs), so that MoE runs see the same total tokens in
roughly twice the number of steps.
The raw-data runs train for significantly more steps due to the
10$\times$ larger corpus, reaching lower final losses (1.69 for dense,
1.70 for MoE) compared to the curated and hybrid runs ($\approx$2.4--3.1).


% -----------------------------------------------------------------------------
\section{Evaluation}
\label{sec:evaluation}
% -----------------------------------------------------------------------------

We evaluate all models using the EleutherAI LM Evaluation Harness, the
same zero-shot benchmark backend used by NeMo Evaluator.  Because our
custom Nemotron-Nano architecture is not natively registered in the NeMo
model zoo, we integrate with the harness directly: a custom wrapper class
(\texttt{NemotronNanoEvalLM}) implements the \texttt{lm\_eval.api.model.LM}
interface, providing batched log-likelihood scoring and greedy generation
on top of our trained checkpoints.  This produces results identical to
those obtained through NeMo Evaluator's harness integration, while
supporting our non-standard model class.

Each checkpoint is evaluated by running:

\begin{verbatim}
python evaluate.py \
  --checkpoint runs/<experiment>/best.pt \
  --tokenizer_dir tokenizer_16k \
  --tasks hellaswag,piqa,arc_easy,arc_challenge,\
         winogrande,sciq,boolq,lambada_openai,wikitext \
  --batch_size 64
\end{verbatim}

\subsection{Benchmarks}

We select a suite of nine zero-shot tasks appropriate for base (non-chat)
language models, covering commonsense reasoning, science, reading
comprehension, and language modeling:

\paragraph{Accuracy benchmarks (higher is better).}
\begin{itemize}
    \item \textbf{HellaSwag}: Commonsense sentence completion (acc\_norm).
    \item \textbf{PIQA}: Physical intuition QA (acc).
    \item \textbf{ARC-Easy / ARC-Challenge}: Grade-school science
          questions (acc\_norm).
    \item \textbf{SciQ}: Multiple-choice science QA (acc\_norm).
    \item \textbf{BoolQ}: Yes/no reading comprehension (acc).
    \item \textbf{WinoGrande}: Coreference resolution (acc).
\end{itemize}

\paragraph{Perplexity benchmarks (lower is better).}
\begin{itemize}
    \item \textbf{WikiText}: Word-level perplexity on WikiText-103.
    \item \textbf{LAMBADA}: Perplexity on the LAMBADA last-word
          prediction task.
\end{itemize}

All evaluations use zero-shot prompting with no task-specific fine-tuning.
Raw evaluation outputs are saved per-run in
\texttt{runs/<experiment>/eval\_results.json}.


% -----------------------------------------------------------------------------
\section{Results}
\label{sec:results}
% -----------------------------------------------------------------------------

\subsection{Group 1: Model Scaling}

Figure~\ref{fig:group1a} shows the effect of scaling model size from 4M to
58M parameters on raw Wikipedia data.  Larger models improve consistently
across all benchmarks, with particularly large gains on SciQ
($+$25.8\,pp from 4M to 58M) and BoolQ ($+$5.0\,pp).

\begin{figure}[h]
\centering
\includegraphics[width=\linewidth]{slides/group1a_model_size.pdf}
\caption{Effect of model size on benchmark accuracy (raw data, dense
architecture).  Scaling from 4M to 58M parameters yields consistent
improvements, confirming that model capacity is a primary performance driver.}
\label{fig:group1a}
\end{figure}

At the fixed 58M model size, Figure~\ref{fig:group1b} compares five data
source configurations.  Performance differences are modest across most
benchmarks, with the notable exception of SciQ where the full raw dataset
(5.7B tokens) dramatically outperforms all subsets ($\approx$328M tokens),
suggesting that data quantity dominates quality effects at this small model
scale.

\begin{figure}[h]
\centering
\includegraphics[width=\linewidth]{slides/group1b_data_source.pdf}
\caption{Effect of data source and quantity at fixed 58M model size.
Most benchmarks show little sensitivity to data source; SciQ is the
exception, where more data helps substantially.}
\label{fig:group1b}
\end{figure}


\subsection{Group 2: Data Quality at Equal Token Budget}

When we control for token count (5.7B tokens each), data quality differences
become clearly visible (Figure~\ref{fig:group2_acc},
Table~\ref{tab:group2_results}).

\begin{figure}[h]
\centering
\includegraphics[width=\linewidth]{slides/group2_equal_accuracy.pdf}
\caption{Benchmark accuracy at equal token budget (5.7B tokens, 382M
dense).  Hybrid data outperforms both raw and curated on most benchmarks.}
\label{fig:group2_acc}
\end{figure}

\begin{table}[h]
\centering
\caption{Group 2 results: accuracy (\%) and perplexity at 5.7B tokens.
Best result per metric in \textbf{bold}.}
\label{tab:group2_results}
\begin{tabular}{l*{7}{c}cc}
\toprule
& \rotatebox{60}{HellaSwag} & \rotatebox{60}{PIQA} & \rotatebox{60}{ARC-E}
& \rotatebox{60}{ARC-C} & \rotatebox{60}{SciQ} & \rotatebox{60}{BoolQ}
& \rotatebox{60}{WinoGrande} & \rotatebox{60}{WikiText$\downarrow$}
& \rotatebox{60}{LAMBADA$\downarrow$} \\
\midrule
Raw     & 31.3 & 61.4 & 45.5 & 24.9 & 61.2 & \textbf{62.1} & \textbf{50.9} & \textbf{36.2} & \textbf{82.0} \\
Curated & 33.5 & 62.3 & 49.9 & 27.2 & 69.6 & 61.3 & 50.5 & 36.4 & 95.8 \\
Hybrid  & \textbf{33.3} & \textbf{64.1} & \textbf{51.0} & \textbf{29.8} & \textbf{67.8} & \textbf{62.1} & 48.6 & 37.1 & 108.4 \\
\bottomrule
\end{tabular}
\end{table}

The hybrid corpus achieves the best accuracy on 5 of 7 benchmarks,
including substantial leads on ARC-Challenge ($+$4.9\,pp over raw) and
ARC-Easy ($+$5.5\,pp over raw).  Raw data retains an edge on perplexity
metrics, likely because its distribution is closer to the WikiText and
LAMBADA test sets.  These results confirm that \emph{at a fixed token
budget, data curation and augmentation meaningfully improve downstream
task performance}.


\subsection{Group 3: Full-Scale Data $\times$ Architecture}

Table~\ref{tab:group3_results} and Figure~\ref{fig:group3_acc} present
the full $3\times2$ grid of results.

\begin{figure}[h]
\centering
\includegraphics[width=\linewidth]{slides/group3_fullscale_accuracy.pdf}
\caption{Full-scale accuracy: Dense vs.\ MoE across three data sources.
MoE consistently outperforms the matched dense model.}
\label{fig:group3_acc}
\end{figure}

\begin{table}[h]
\centering
\caption{Group 3 results: full-scale runs.  Best per data source in
\textbf{bold}; overall best \underline{underlined}.}
\label{tab:group3_results}
\resizebox{\linewidth}{!}{%
\begin{tabular}{ll*{7}{c}cc}
\toprule
\textbf{Data} & \textbf{Arch.}
& \rotatebox{60}{HellaSwag} & \rotatebox{60}{PIQA} & \rotatebox{60}{ARC-E}
& \rotatebox{60}{ARC-C} & \rotatebox{60}{SciQ} & \rotatebox{60}{BoolQ}
& \rotatebox{60}{WinoGrande} & \rotatebox{60}{WikiText$\downarrow$}
& \rotatebox{60}{LAMBADA$\downarrow$} \\
\midrule
Raw     & Dense & 43.0 & 68.0 & 55.2 & 29.2 & 76.4 & 60.0 & 51.8 & 20.9 & 21.2 \\
Raw     & MoE   & \underline{\textbf{49.1}} & \underline{\textbf{70.2}} & \underline{\textbf{59.6}} & \underline{\textbf{32.1}} & \underline{\textbf{79.6}} & \textbf{60.7} & \underline{\textbf{54.5}} & \underline{\textbf{17.6}} & \underline{\textbf{15.2}} \\
\midrule
Curated & Dense & 33.5 & 62.3 & 49.9 & 27.2 & 69.6 & \textbf{61.3} & \textbf{50.5} & 36.4 & 95.8 \\
Curated & MoE   & \textbf{36.5} & \textbf{64.6} & \textbf{52.1} & \textbf{28.9} & 68.9 & 59.6 & 49.6 & \textbf{30.5} & \textbf{59.9} \\
\midrule
Hybrid  & Dense & 34.7 & 63.7 & 53.1 & 29.0 & 70.4 & 61.3 & 49.5 & 35.9 & 85.1 \\
Hybrid  & MoE   & \textbf{38.5} & \textbf{66.2} & \textbf{57.6} & \textbf{31.7} & \textbf{73.0} & \textbf{61.4} & \textbf{52.0} & \textbf{28.6} & \textbf{55.8} \\
\bottomrule
\end{tabular}%
}
\end{table}

\paragraph{Finding 1: Volume beats quality at scale.}
The raw-data models dominate the overall leaderboard despite using
unfiltered data.  With $\approx$55B training tokens (10$\times$ the
curated/hybrid corpora), the raw MoE achieves the best score on every
benchmark.  HellaSwag reaches 49.1\% (vs.\ 38.5\% for hybrid MoE),
and WikiText perplexity drops to 17.6 (vs.\ 28.6).  At current corpus
sizes, we do not yet have enough curated data for quality to overcome
the sheer advantage of quantity.

\paragraph{Finding 2: Curation helps per-token.}
Comparing the curated and hybrid columns against raw \emph{at similar
token counts} (5--9B), hybrid data consistently outperforms raw:
hybrid MoE beats curated MoE on 6 of 7 accuracy benchmarks and
achieves 7.3 points lower WikiText perplexity than the corresponding
dense model.  This confirms that data curation and synthetic
augmentation improve sample efficiency.

\paragraph{Finding 3: MoE $>$ Dense for free.}
Across all three data sources, the MoE architecture improves over the
matched dense model on nearly every metric, despite using the same
number of active parameters per token.  The gains are:
\begin{itemize}
    \item HellaSwag: $+$3.0 to $+$6.1\,pp
    \item PIQA: $+$2.2 to $+$2.4\,pp
    \item ARC-Easy: $+$2.3 to $+$4.5\,pp
    \item WikiText PPL: $-$3.3 to $-$7.3
\end{itemize}
The MoE advantage is largest on the raw corpus (where more training
data allows the experts to specialize more effectively) and smallest
on the curated corpus.

\begin{figure}[h]
\centering
\includegraphics[width=0.85\linewidth]{slides/group3_heatmap.pdf}
\caption{Heatmap of the $3\times2$ experimental grid.  Each cell shows
accuracy (\%) for a given benchmark; the rightmost column shows the
average across all seven benchmarks.  MoE consistently outperforms
Dense within each data source.}
\label{fig:group3_heatmap}
\end{figure}

Figure~\ref{fig:group3_ppl} shows the perplexity comparison.  The
pattern mirrors the accuracy results: raw-data models achieve the
lowest perplexity, and MoE improves over dense in every case.

\begin{figure}[h]
\centering
\includegraphics[width=\linewidth]{slides/group3_fullscale_perplexity.pdf}
\caption{WikiText and LAMBADA perplexity for the six full-scale runs.
Lower is better.  MoE reduces perplexity across all data sources.}
\label{fig:group3_ppl}
\end{figure}

\subsection{Summary}

Our experiments reveal a clear hierarchy: (1)~data volume is the
strongest lever at current scales, (2)~data curation improves
sample efficiency when token budgets are matched, and (3)~sparse
MoE routing provides a consistent ``free lunch'' over dense models
at equivalent per-token compute.  The best overall model---raw-data
MoE with 1.4B total parameters and 390M active---achieves 49.1\%
HellaSwag, 70.2\% PIQA, and 17.6 WikiText perplexity.
